
\documentclass{article}
\usepackage{colortbl}
\usepackage{makecell}
\usepackage{multirow}
\usepackage{supertabular}

\begin{document}

\newcounter{utterance}

\centering \large Interaction Transcript for game `cladder', experiment `full\_v1.5\_default', episode 7911 with score{-}Qwen3.5{-}9B{-}sp.
\vspace{24pt}

{ \footnotesize  \setcounter{utterance}{1}
\setlength{\tabcolsep}{0pt}
\begin{supertabular}{c@{$\;$}|p{.15\linewidth}@{}p{.15\linewidth}p{.15\linewidth}p{.15\linewidth}p{.15\linewidth}p{.15\linewidth}}
    \# & \multicolumn{2}{c}{Player} && \multicolumn{2}{c}{Game Master} \\
    \hline

    \theutterance \stepcounter{utterance}  
    & & & \multicolumn{4}{p{0.6\linewidth}}{
        \cellcolor[rgb]{0.9,0.9,0.9}{
            \makecell[{{p{\linewidth}}}]{
                \texttt{\tiny{[P1$\langle$GM]}}
                \texttt{You are an expert in causal inference. The following question is not a typical commonsense query, but rather a meticulously designed question created by a professor specializing in causal inference, intended to assess the students' mastery of the course content.} \\
\texttt{Imagine a self{-}contained, hypothetical world with only the following conditions, and without any unmentioned factors or causal relationships: Rixq has a direct effect on zuph and xevu. Swoy has a direct effect on zuph and xevu. Zuph has a direct effect on xevu. Swoy is unobserved. The overall probability of rixq is 3\%. The probability of not rixq and xevu is 78\%. The probability of rixq and xevu is 2\%. Is the chance of xevu larger when observing rixq?} \\
\texttt{Start your answer with 'Yes' or 'No', followed by additional reasoning or evidence to support your explanation.} \\
            }
        }
    }
     \\ \\

    \theutterance \stepcounter{utterance}  
    & \multicolumn{4}{p{0.6\linewidth}}{
        \cellcolor[rgb]{0.561,0.737,0.561}{
            \makecell[{{p{\linewidth}}}]{
                \texttt{\tiny{[P1$\rangle$GM]}}
                \texttt{No} \\
\\ 
\texttt{**Reasoning:**} \\
\\ 
\texttt{To determine if the chance of \$xevu\$ is larger when observing \$rixq\$, we must compare the conditional probability \$P(xevu | rixq)\$ with the marginal probability \$P(xevu)\$.} \\
\\ 
\texttt{1.  **Calculate the marginal probability \$P(xevu)\$:**} \\
\texttt{    We are given the joint probabilities involving \$rixq\$ and \$xevu\$:} \\
\texttt{    *   \$P(\textbackslash{}neg rixq \textbackslash{}cap xevu) = 78\textbackslash{}\%\$} \\
\texttt{    *   \$P(rixq \textbackslash{}cap xevu) = 2\textbackslash{}\%\$} \\
\\ 
\texttt{    The total probability of \$xevu\$ is the sum of these two mutually exclusive cases:} \\
\texttt{    \$\$P(xevu) = P(\textbackslash{}neg rixq \textbackslash{}cap xevu) + P(rixq \textbackslash{}cap xevu) = 78\textbackslash{}\% + 2\textbackslash{}\% = 80\textbackslash{}\%\$\$} \\
\\ 
\texttt{2.  **Calculate the conditional probability \$P(xevu | rixq)\$:**} \\
\texttt{    By definition, \$P(xevu | rixq) = \textbackslash{}frac\{P(rixq \textbackslash{}cap xevu)\}\{P(rixq)\}\$.} \\
\texttt{    *   We are given \$P(rixq) = 3\textbackslash{}\%\$.} \\
\texttt{    *   We are given \$P(rixq \textbackslash{}cap xevu) =} \\
            }
        }
    }
    & & \\ \\

    \theutterance \stepcounter{utterance}  
    & & & \multicolumn{2}{p{0.3\linewidth}}{
        \cellcolor[rgb]{0.9,0.9,0.9}{
            \makecell[{{p{\linewidth}}}]{
                \texttt{\tiny{[GM$|$GM]}}
                \texttt{No} \\
            }
        }
    }
    & & \\ \\

    \theutterance \stepcounter{utterance}  
    & & & \multicolumn{2}{p{0.3\linewidth}}{
        \cellcolor[rgb]{0.9,0.9,0.9}{
            \makecell[{{p{\linewidth}}}]{
                \texttt{\tiny{[GM$|$GM]}}
                \texttt{no} \\
            }
        }
    }
    & & \\ \\

    \theutterance \stepcounter{utterance}  
    & & & \multicolumn{2}{p{0.3\linewidth}}{
        \cellcolor[rgb]{0.9,0.9,0.9}{
            \makecell[{{p{\linewidth}}}]{
                \texttt{\tiny{[GM$|$GM]}}
                \texttt{game\_result = WIN} \\
            }
        }
    }
    & & \\ \\

\end{supertabular}
}

\end{document}
